\documentclass[11pt]{article}

% Shared notes settings for Elements of AIML lectures (UPES).
% Keep this file free of \documentclass to allow reuse via \input.

\usepackage[utf8]{inputenc}
\usepackage[T1]{fontenc}
\usepackage[a4paper,margin=1in]{geometry}
\usepackage{hyperref}
\usepackage{xurl}
\usepackage{booktabs}
\usepackage{amsmath}
\usepackage{amssymb}
\usepackage{listings}
\usepackage{xcolor}
\usepackage{enumitem}
\usepackage{parskip}

% Code listings (general-purpose).
\lstset{
  basicstyle=\ttfamily\small,
  keywordstyle=\color{blue},
  commentstyle=\color{gray},
  stringstyle=\color{teal},
  showstringspaces=false,
  columns=fullflexible,
  frame=single,
  framerule=0.2pt,
  breaklines=true
}

\setlist[itemize]{topsep=0.3em,itemsep=0.2em}
\setlist[enumerate]{topsep=0.3em,itemsep=0.2em}

% Simple per-section source line.
\newcommand{\notesources}[1]{\par\noindent{\small\color{gray}\textit{Sources:} \sloppy #1\par}}

% Unit-wise source strings for this subject (used by slides/notes).
% Path is relative to the lecture `latex/` folder (the compilation CWD).
\IfFileExists{../../../_shared/aiml-sources.tex}{%
  % Source strings for Elements of AIML (used by slides + notes).

\newcommand{\AIMLRepoURL}{https://github.com/tali7c/Elements-of-AIML}

\newcommand{\AIMLLabSources}{%
Provost and Fawcett, \textit{Data Science for Business}.;%
McKinney, \textit{Python for Data Analysis}.;%
WEKA Documentation (Waikato Environment for Knowledge Analysis), accessed 2026-02-28.;%
Matplotlib and Seaborn documentation, accessed 2026-02-28.%
}

\newcommand{\AIMLUnitISources}{%
Rich and Knight, \textit{Artificial Intelligence}, McGraw Hill, 2017.;%
Russell and Norvig, \textit{Artificial Intelligence: A Modern Approach} (4e), Ch. 1--2 (Introduction, Intelligent Agents).;%
Poole and Mackworth, \textit{Artificial Intelligence: Foundations of Computational Agents} (2e), selected sections.%
}

\newcommand{\AIMLUnitIISources}{%
Russell and Norvig, \textit{Artificial Intelligence: A Modern Approach} (4e), Ch. 7--9 (Logical Agents, First-Order Logic, Inference).;%
Huth and Ryan, \textit{Logic in Computer Science} (2e), selected sections on propositional and first-order logic.;%
Genesereth and Nilsson, \textit{Logical Foundations of Artificial Intelligence}, selected chapters.%
}

\newcommand{\AIMLUnitIIISources}{%
Mueller and Massaron, \textit{Machine Learning for Dummies}, For Dummies, 2016.;%
Mitchell, \textit{Machine Learning}, selected chapters on learning basics and evaluation.;%
Scikit-learn documentation, accessed 2026-02-28.%
}

\newcommand{\AIMLUnitIVSources}{%
Mueller and Massaron, \textit{Machine Learning for Dummies}, For Dummies, 2016.;%
James, Witten, Hastie, and Tibshirani, \textit{An Introduction to Statistical Learning} (2e), selected chapters.;%
Scikit-learn documentation, accessed 2026-02-28.%
}

\newcommand{\AIMLUnitVSources}{%
NIST, \textit{AI Risk Management Framework (AI RMF 1.0)}, 2023.;%
Barocas, Hardt, and Narayanan, \textit{Fairness and Machine Learning} (online book), selected sections.;%
Knaflic, \textit{Storytelling with Data}, selected chapters on communicating results.%
}
%
}{%
  % Faculty copies live under `private/` so their compile CWD is different.
  \IfFileExists{../../../../public/_shared/aiml-sources.tex}{%
    % Source strings for Elements of AIML (used by slides + notes).

\newcommand{\AIMLRepoURL}{https://github.com/tali7c/Elements-of-AIML}

\newcommand{\AIMLLabSources}{%
Provost and Fawcett, \textit{Data Science for Business}.;%
McKinney, \textit{Python for Data Analysis}.;%
WEKA Documentation (Waikato Environment for Knowledge Analysis), accessed 2026-02-28.;%
Matplotlib and Seaborn documentation, accessed 2026-02-28.%
}

\newcommand{\AIMLUnitISources}{%
Rich and Knight, \textit{Artificial Intelligence}, McGraw Hill, 2017.;%
Russell and Norvig, \textit{Artificial Intelligence: A Modern Approach} (4e), Ch. 1--2 (Introduction, Intelligent Agents).;%
Poole and Mackworth, \textit{Artificial Intelligence: Foundations of Computational Agents} (2e), selected sections.%
}

\newcommand{\AIMLUnitIISources}{%
Russell and Norvig, \textit{Artificial Intelligence: A Modern Approach} (4e), Ch. 7--9 (Logical Agents, First-Order Logic, Inference).;%
Huth and Ryan, \textit{Logic in Computer Science} (2e), selected sections on propositional and first-order logic.;%
Genesereth and Nilsson, \textit{Logical Foundations of Artificial Intelligence}, selected chapters.%
}

\newcommand{\AIMLUnitIIISources}{%
Mueller and Massaron, \textit{Machine Learning for Dummies}, For Dummies, 2016.;%
Mitchell, \textit{Machine Learning}, selected chapters on learning basics and evaluation.;%
Scikit-learn documentation, accessed 2026-02-28.%
}

\newcommand{\AIMLUnitIVSources}{%
Mueller and Massaron, \textit{Machine Learning for Dummies}, For Dummies, 2016.;%
James, Witten, Hastie, and Tibshirani, \textit{An Introduction to Statistical Learning} (2e), selected chapters.;%
Scikit-learn documentation, accessed 2026-02-28.%
}

\newcommand{\AIMLUnitVSources}{%
NIST, \textit{AI Risk Management Framework (AI RMF 1.0)}, 2023.;%
Barocas, Hardt, and Narayanan, \textit{Fairness and Machine Learning} (online book), selected sections.;%
Knaflic, \textit{Storytelling with Data}, selected chapters on communicating results.%
}
%
  }{%
    % Source strings for Elements of AIML (used by slides + notes).

\newcommand{\AIMLRepoURL}{https://github.com/tali7c/Elements-of-AIML}

\newcommand{\AIMLLabSources}{%
Provost and Fawcett, \textit{Data Science for Business}.;%
McKinney, \textit{Python for Data Analysis}.;%
WEKA Documentation (Waikato Environment for Knowledge Analysis), accessed 2026-02-28.;%
Matplotlib and Seaborn documentation, accessed 2026-02-28.%
}

\newcommand{\AIMLUnitISources}{%
Rich and Knight, \textit{Artificial Intelligence}, McGraw Hill, 2017.;%
Russell and Norvig, \textit{Artificial Intelligence: A Modern Approach} (4e), Ch. 1--2 (Introduction, Intelligent Agents).;%
Poole and Mackworth, \textit{Artificial Intelligence: Foundations of Computational Agents} (2e), selected sections.%
}

\newcommand{\AIMLUnitIISources}{%
Russell and Norvig, \textit{Artificial Intelligence: A Modern Approach} (4e), Ch. 7--9 (Logical Agents, First-Order Logic, Inference).;%
Huth and Ryan, \textit{Logic in Computer Science} (2e), selected sections on propositional and first-order logic.;%
Genesereth and Nilsson, \textit{Logical Foundations of Artificial Intelligence}, selected chapters.%
}

\newcommand{\AIMLUnitIIISources}{%
Mueller and Massaron, \textit{Machine Learning for Dummies}, For Dummies, 2016.;%
Mitchell, \textit{Machine Learning}, selected chapters on learning basics and evaluation.;%
Scikit-learn documentation, accessed 2026-02-28.%
}

\newcommand{\AIMLUnitIVSources}{%
Mueller and Massaron, \textit{Machine Learning for Dummies}, For Dummies, 2016.;%
James, Witten, Hastie, and Tibshirani, \textit{An Introduction to Statistical Learning} (2e), selected chapters.;%
Scikit-learn documentation, accessed 2026-02-28.%
}

\newcommand{\AIMLUnitVSources}{%
NIST, \textit{AI Risk Management Framework (AI RMF 1.0)}, 2023.;%
Barocas, Hardt, and Narayanan, \textit{Fairness and Machine Learning} (online book), selected sections.;%
Knaflic, \textit{Storytelling with Data}, selected chapters on communicating results.%
}
%
  }%
}


\title{Elements of AIML\\Unit 04 -- Lecture 01 Notes\\Supervised Learning (Regression)}
\author{Tofik Ali}
\date{\today}

\begin{document}
\maketitle
\tableofcontents

\section{Lecture Overview}
This lecture introduces \textbf{regression}, a supervised learning task where the target is a continuous value.
We discuss linear regression, loss functions, evaluation metrics, and overfitting/regularization.
\notesources{\AIMLUnitIVSources}

\section{How to use these notes (self-study)}
For regression, always separate three things: (1) the model form ($\hat{y}=f(x)$), (2) the loss used for training (e.g., MSE), and (3) the metric used for evaluation (MAE/RMSE/$R^2$).
Work through the mini metric computation example and then attempt the practice questions.

\section{Learning Outcomes}
\begin{itemize}[leftmargin=*]
  \item Define regression and give real-world examples.
  \item Explain linear regression and the role of a loss function.
  \item Compute and interpret MAE, RMSE, and $R^2$.
  \item Explain overfitting and the purpose of regularization.
\end{itemize}

\section{Keywords}
\begin{itemize}[leftmargin=*]
  \item \textbf{Regression:} predicting a continuous target.
  \item \textbf{Loss function:} objective minimized during training (e.g., MSE).
  \item \textbf{MAE/RMSE:} error metrics in target units.
  \item \textbf{$R^2$:} variance-explained metric (relative measure).
  \item \textbf{Overfitting:} fits training noise; poor generalization.
  \item \textbf{Regularization:} penalty term to control complexity.
\end{itemize}

\section{Abbreviations and symbols}
\begin{itemize}[leftmargin=*]
  \item $X$: feature matrix, $y$: true targets, $\hat{y}$: predictions.
  \item $n$: number of training examples, $d$: number of features.
  \item $w$: weight vector (parameters), $w_0$: intercept (bias term).
  \item $e_i = y_i - \hat{y}_i$: residual/error for example $i$.
  \item MAE/MSE/RMSE: regression metrics; $R^2$: coefficient of determination.
  \item $\lambda$: regularization strength (Ridge/Lasso).
\end{itemize}

\section{Regression basics}
In regression, we learn a function $f$ so that:
\[
  \hat{y} = f(x)
\]
where $y$ is continuous (e.g., price). Regression is used in forecasting, estimation, and prediction problems.

\subsection{Regression vs classification (quick reminder)}
\begin{itemize}[leftmargin=*]
  \item \textbf{Regression}: output is a real number (temperature, price).
  \item \textbf{Classification}: output is a class label (spam/ham).
\end{itemize}

\section{Linear regression}
\subsection{Model}
For $d$ features:
\[
  \hat{y} = w_0 + \sum_{j=1}^{d} w_j x_j
\]
The goal is to find weights $w$ that minimize prediction error on training data.

\subsection{Vector/matrix form (compact notation)}
If we add a constant ``1'' feature for the intercept, we can write:
\[
  \hat{y} = X w
\]
where $X \in \mathbb{R}^{n \times (d+1)}$ and $w \in \mathbb{R}^{(d+1)}$.
This form is useful for reasoning about optimization and implementation.

\subsection{Loss functions}
A common loss is mean squared error (MSE):
\[
  \text{MSE} = \frac{1}{n}\sum_{i=1}^{n}(y_i - \hat{y}_i)^2
\]
RMSE is $\sqrt{\text{MSE}}$ and is easier to interpret because it is in the same units as $y$.
MAE uses absolute values:
\[
  \text{MAE} = \frac{1}{n}\sum_{i=1}^{n}|y_i - \hat{y}_i|
\]

\subsection{MAE vs MSE (why do we have both?)}
\begin{itemize}[leftmargin=*]
  \item MSE/RMSE penalizes large errors strongly (squares the error), so it is sensitive to outliers.
  \item MAE penalizes errors linearly, often more robust to outliers.
  \item In practice, choose based on the cost of large mistakes (domain-dependent).
\end{itemize}

\section{Evaluation metrics}
\subsection{MAE vs RMSE}
\begin{itemize}[leftmargin=*]
  \item MAE treats all errors linearly.
  \item RMSE penalizes large errors more strongly.
\end{itemize}
Choice depends on whether large errors are especially costly.

\subsection{$R^2$ (coefficient of determination)}
One common definition:
\[
  R^2 = 1 - \frac{\sum_i (y_i-\hat{y}_i)^2}{\sum_i (y_i-\bar{y})^2}
\]
$R^2$ is relative to a baseline that predicts the mean $\bar{y}$.
It can be misleading when:
\begin{itemize}[leftmargin=*]
  \item comparing across datasets with different variance,
  \item evaluating non-linear relationships with a poor model,
  \item using it without checking residual patterns.
\end{itemize}

\subsection{Mini worked example (compute MAE/RMSE)}
Suppose:
\[
  y=[50,60,80],\quad \hat{y}=[55,58,70]
\]
Residuals: $e = y-\hat{y}=[-5,2,10]$.
\begin{itemize}[leftmargin=*]
  \item MAE $= \frac{|{-5}|+|2|+|10|}{3} = \frac{17}{3} \approx 5.67$
  \item MSE $= \frac{25+4+100}{3} = 43$, so RMSE $=\sqrt{43}\approx 6.56$
\end{itemize}

\section{Underfitting, overfitting, and regularization}
\subsection{Underfitting vs overfitting}
\begin{itemize}[leftmargin=*]
  \item Underfitting: too simple; high error on train and test.
  \item Overfitting: too complex; low train error, high test error.
\end{itemize}
Always evaluate using proper splits/cross-validation (Unit III).

\subsection{Residual analysis (basic sanity check)}
Even if metrics look good, we should inspect residuals:
\begin{itemize}[leftmargin=*]
  \item If residuals show patterns (e.g., curve shape), a linear model may be misspecified.
  \item If residual spread increases with $x$, variance is not constant (heteroscedasticity).
  \item Large residuals may indicate outliers or missing features.
\end{itemize}

\subsection{Regularization idea}
Regularization adds a penalty to control complexity. Example (Ridge):
\[
  \text{Loss} = \text{MSE} + \lambda \|w\|^2
\]
Lasso uses $\|w\|_1$ and can push weights to zero.

\subsection{Ridge vs Lasso (high-level)}
\begin{itemize}[leftmargin=*]
  \item \textbf{Ridge (L2)}: shrinks weights smoothly; helps with multicollinearity.
  \item \textbf{Lasso (L1)}: can set some weights to exactly zero; acts like feature selection.
\end{itemize}
Regularization reduces variance (overfitting) at the cost of some bias.

\section{Minimal Python example (illustrative)}
\begin{lstlisting}[language=Python]
from sklearn.model_selection import train_test_split
from sklearn.linear_model import LinearRegression, Ridge
from sklearn.metrics import mean_absolute_error, mean_squared_error
import numpy as np

X_train, X_test, y_train, y_test = train_test_split(X, y, test_size=0.2, random_state=42)
model = LinearRegression().fit(X_train, y_train)
pred = model.predict(X_test)

mae = mean_absolute_error(y_test, pred)
rmse = np.sqrt(mean_squared_error(y_test, pred))
print(mae, rmse)
\end{lstlisting}

\section{Common pitfalls (and how to avoid them)}
\begin{itemize}[leftmargin=*]
  \item \textbf{Evaluating on training data:} this hides overfitting; always use proper splits/CV.
  \item \textbf{Ignoring units:} RMSE is in the same units as $y$; MSE is in squared units.
  \item \textbf{Over-trusting $R^2$:} check residuals and absolute error; $R^2$ can be misleading across datasets.
  \item \textbf{Forgetting preprocessing:} outliers, scaling, and feature encoding can dominate performance.
\end{itemize}

\section{Practice (with answers)}
\subsection*{P1. When would you prefer MAE over RMSE?}
\textbf{Answer:} when you want robustness to outliers or when large errors should not dominate the metric too strongly.

\subsection*{P2. What does $\lambda$ do in Ridge regression?}
\textbf{Answer:} it controls the penalty strength; larger $\lambda$ shrinks weights more (simpler model), often reducing overfitting.

\subsection*{P3. Give one reason why $R^2$ can be misleading.}
\textbf{Answer:} it depends on variance of $y$; a model can have a high $R^2$ but still have large errors in absolute terms, or a low $R^2$ even when RMSE is acceptable.

\section{Exit questions (with answers)}

\subsection*{Q1. Regression vs classification?}
\textbf{Answer:} regression predicts a continuous value; classification predicts a discrete class label.

\subsection*{Q2. Why is RMSE easier to interpret than MSE?}
\textbf{Answer:} RMSE has the same units as the target variable (MSE is squared units).

\subsection*{Q3. What does $R^2$ measure and when is it misleading?}
\textbf{Answer:} $R^2$ measures variance explained relative to predicting the mean; it can be misleading across datasets, with non-linear patterns, or without residual analysis.

\subsection*{Q4. Overfitting mitigation strategies?}
\textbf{Answer:} regularization, more data, simpler model, better features, cross-validation, early stopping (for iterative models).

\section{Summary}
Regression is a core supervised learning task. Correct evaluation and controlling overfitting are essential for reliable predictions.
\notesources{\AIMLUnitIVSources}

\end{document}
